\newpage
\pagestyle{plain}

\section{\centering ЗАКЛЮЧЕНИЕ}

%цели, выводы, заключение - краткий пересказ, зеркально к введению

В настоящей работе было проведено исследование способов повышения точности численного решения систем обыкновенных дифференциальных уравнений, решения которых содержат осцилляции. Были разработаны и исследованы модификации параметрических схем Рунге--Кутты, позволяющие адаптировать коэффициенты под каждую из поставленных задач, а также проведены дополнительные эксперименты для демонстрации свойств изучаемых схем.

В ходе выполнения работы были решены следующие задачи:
\begin{enumerate}
	\item проведён теоретический анализ методов Рунге--Кутты 6-го порядка и принципов построения их параметрических семейств;
	\item задана математическая модель оптимизации свободных параметров и выбрана целевая функция, минимизирующая локальную ошибку;
	\item на языке Python реализован и протестирован на четырёх тестовых задачах модульный программный комплекс;
	\item выполнены численные эксперименты, визуализация полученных решений и проведено сравнение результатов разработанного метода с результатами источников.
	\item визуализированы и исследованы свойства схем, такие как скорость сходимости методов, сохранение методами теоретически инвариантных величин, практический порядок точности и возрастание времени исполнения схемы в зависимости от размерности системы ОДУ.
\end{enumerate}

Таким образом, разработанный параметрический метод преодолевает основные сложности, возникающие при решении задач с осциллирующими решениями, и подтверждает свою эффективность на тестовых задачах математики и физики.

Перечислим основные из возможных направлений дальнейших исследований:

\begin{enumerate}
	\item исследовать влияние альтернативных норм (например, первой нормы $L_1$ или Чебышёвской нормы) в критерии качества на скорость сходимости алгоритмов и их точность;
	\item провести анализ алгоритма на других реальных задачах, не имеющих точного аналитического решения;
	\item провести анализ стабильности инвариантов движения, увеличив временные интервалы;
	\item использовать параллельные вычисления при расчёте значений функционала ошибки и оценить эффективность параллелизации;
	\item изучить влияние шага дискретизации на свойства детерминированного сеточного поиска минимума функционала;
	\item провести сравнение эффективности метода роя частиц и других безградиентных алгоритмов.
\end{enumerate}

Перспективным направлением в дальнейшем развитии построенного вычислительного комплекса станет интеграция разработанных алгоритмов с методами машинного обучения. В частности, применение искусственных нейронных сетей позволит аппроксимировать сложный многомодальный рельеф целевых функций и на основе предобучения предсказывать субоптимальные значения вектора параметров. Такой подход обеспечит возможность масштабирования параметрических схем на системы большей размерности без высоких вычислительных затрат.